\mnImportant
\PID
Мерење физичке величине $x$ обавља се мерним инструментом са случајном несигурношћу мерења 
$u_x$. Мерење се понваља $n$ пута, чиме се добија средња вредност мерења 
$$
{\overline x = 
\dfrac{x_1 + x_2 + \cdots + x_n}{n}}.
$$ Одредити израз за случнајну мерну несигурност, $u_{\overline x}$ добијене средње 
вредности. 

\RESENJE 

Случајна несигурност величине дате изразом $y = y(x_1, x_2, \ldots x_n)$, у случају \textit{некорелисаних} мерења, може се 
записати у облику\footnote{Разлог за забирање \textit{квадрата} случајних грешака може се пронаћи у чињеници да се 
\textit{снаге} некорелисаних сигнала сабирају, односно важи 
$ \overline{(x+y)^2} = \overline{x^2} + \overline{y^2} + 2\cancelto{0}{\overline{xy}}$, пошто претпоставка 
о некорелисаности значи да је управо средња вредност производа једнака нули. Из потпуно истог разлога се 
сабирају \textit{снаге} извора некорелисаних шумова, што би читаоцу могло да буде познато из области 
аналогне електронике, где се записује $e_n^2 = e_1^2 + e_2^2 + \cdots + e_n^2$. }
\begin{equation}
    u_y = \sqrt{  \sum_{i = 1}^n \left( \dfrac{\partial y}{\partial x_i} u_{x_i}  \right)^2  }.
    \label{eq:greska_def}
\end{equation} 
У разматраном случају, је $y = \overline x$, па је због симетрије, 
$\dfrac{\partial y}{\partial i} = \dfrac{1}{n}$, док је $u_{x_i} = u_x$ па се заменом у 
\eqref{eq:greska_def} добија
\begin{equation}
    u_{\overline x} = 
    \sqrt{
        \sum_{i = 1}^n \left( \dfrac{1}{n} u_x  \right)^2 
    }
    = \dfrac{1}{\sqrt n} \, u_x.
\end{equation}
Дакле, закључујемо да случајна несигурност средње вредност опада са кореном броја поновљених мерења. Овај резутлат 
је основни разлог за усредњавање већег броја мерења. Начелно, ако се мерење понови 100 пута, несигурност средње вредности је 
10 пута мања од несигурности појединачног мерења. Важно је нагласити да је ово статистичка мера, док је апсолутна 
грешка средње вредности и даље непромењена, јер је 
$\Delta\left( \dfrac{\sum_{i=1}^n x_i}{n} \right) = 
\dfrac{\cancel{n} \Delta x}{\cancel{n}}= \Delta x$. Односно, у најгорем случају грешка се усредњавањем не поправља, али 
се статистичка мера несигурности смањује.

